\documentclass[11pt,letterpaper]{article}
\usepackage[margin=1in]{geometry}
\usepackage{fontspec}
\usepackage{xeCJK}
\setmainfont{Times New Roman}
\setCJKmainfont[AutoFakeBold=true]{DFKai-SB}
\usepackage{setspace}
\usepackage{enumitem}
\usepackage[para]{footmisc}
\renewcommand{\footnotesize}{\fontsize{8}{9.2}\selectfont}
\setlength{\skip\footins}{0.5em plus 0.1em minus 0.1em}
\usepackage{hyperref}
\hypersetup{colorlinks=true, allcolors=black}
\setlength{\parindent}{0pt}
\setlength{\parskip}{0.38em}
\setstretch{1.0}
\pagestyle{empty}

\begin{document}

\begin{center}
{\Large 研究志趣與未來研究計畫}\\[0.25em]
許可證（Ke Cheng Hsu）\\[0.15em]
{\small 應徵中華經濟研究院助研究員}
\end{center}

\noindent\textbf{研究動機。}線上市場的運作依賴關於「反悔」的規則：消費者可以退回已購買的商品，書店可以退回賣不掉的存貨，平台可以在交易未完成時退還費用。誰制定這些規則、規則該如何設計，決定了什麼商品被買走、什麼商品被浪費，以及成本由誰承擔。我以實證方法研究這些問題。

\medskip
\noindent\textbf{研究志趣。}我的研究領域為實證產業組織與數位經濟。我研究在不確定性消除之前就必須做出、且事後調整成本高昂的決策。消費者面對產品價值的不確定性，可以在付款前透過搜尋、或在付款後透過退貨來了解產品；賣家面對需求的不確定性，必須事先訂定價格與庫存。一項看似修正某一環節的政策，往往扭曲另一環節，而兩個環節是同時決定的。這種共同決定的關係，貫穿我在產品退貨、消費者搜尋與需求不確定下庫存決策的研究，也是我帶入政策分析的角度。

\medskip
\noindent\textbf{工作論文。}
\begin{itemize}[leftmargin=1.6em,itemsep=0.2em,topsep=0.1em,parsep=0.1em]
\item \textbf{Designing Return Policy in E-Commerce: A Trade-off between Purchase and Return Efficiency}（求職論文）。退貨政策設定兩個邊際：消費者是否購買，以及收到商品後是否退貨。較嚴格的政策抑制購買後無效率的退貨，但退貨減少後，賣家每筆交易保留的金額提高，提高價格變得更有利可圖，進而可能排擠原本有價值的購買。我使用淘寶一家女裝零售商的交易資料，在六個服裝品類、獨占定價下估計購買、保險與退貨決策的結構模型。補貼退貨運費的寬鬆政策在毛利與殘值較高時表現較佳，因為它擴大有效率的購買；較嚴格的政策在毛利較薄時表現較佳，因為它抑制過度退貨。折扣價搭配適度補貨費的選單，能平衡購買效率與退貨效率，在多數品類中優於前兩者。台灣使這個問題具體化：《消費者保護法》第19條賦予網購消費者七日猶豫期，可無理由、無費用退貨。\footnote{消費者保護法第19條。英文條文：全國法規資料庫，\url{https://law.moj.gov.tw/ENG/LawClass/LawAll.aspx?pcode=J0170001}。}在免費退貨成為法定預設之下，賣家在此規範下應如何定價與設計退貨機制，是受規範市場中的實際問題，而非假設性問題。
\item \textbf{Product Design with Heterogeneous Search Costs}（與Pak Hung Au合著）。我們在序貫搜尋模型中引入搜尋成本異質性。低搜尋成本的消費者尋找利基產品，高搜尋成本的消費者購買標準化產品。當搜尋成本下降，更多消費者轉向搜尋利基產品，而這種「誰在搜尋」的轉變可能使兩類產品的價格同時上升。
\end{itemize}

\medskip
\noindent\textbf{進行中的研究。}
\begin{itemize}[leftmargin=1.6em,itemsep=0.2em,topsep=0.1em,parsep=0.1em]
\item \textbf{Pass-Through and Incidence with Product Returns}（與Ce Bian、Liang Zhao合著）。本計畫研究在購買可以反悔時，租稅如何轉嫁。只有當買方保留商品時，賣家才真正收到價款，因此每筆銷售都承擔生產成本、退貨處理成本，以及退貨時無法回收的租稅或費用。價格也會影響退貨率，進而回饋到利潤率與供給量，因此轉嫁可能更強、更弱，甚至為負。忽略退貨邊際的計算，可能使預測的價格與福利效果完全相反。本計畫也探討留下的銷售與退貨應如何課稅，因為實務上銷售稅在退貨時全額退還，平台佣金卻只部分退還。台灣的電商平台市場使這個問題具體化：蝦皮分階段調漲賣家費用，合計抽成據報導已超過14\%，最近的調整包括2026年調高金流與系統處理費，並要求所有賣家加入免運專案；同時，消費者在猶豫期內可無費用退貨，平台對退貨訂單按比例退還成交手續費。\footnote{蝦皮購物賣家幫助中心：成交手續費與金流費，\url{https://seller.shopee.tw/edu/article/26470}；2026年1月費率更新與免運專案，\url{https://seller.shopee.tw/edu/article/26430}；退貨退款訂單不收成交手續費、費用按比例退還，\url{https://seller.shopee.tw/edu/article/12157}。合計抽成超過14\%與賣家反彈：《天下雜誌》專訪蝦皮台灣總經理，2026年，\url{https://www.cw.com.tw/article/5141245}。}這些費用如何轉嫁到價格、以及各項費用在退貨時的退還方式如何影響轉嫁，是台灣平台與賣家之間正在進行的爭議。
\item \textbf{Too Much Unsold Stock?: Environmental Consequences of Price Commitment}（與Kohei Kawaguchi、Kei Kawai合著）。在日本書籍市場，出版商在維持轉售價格下訂定零售價，書店將未售出的書籍退回銷毀。固定零售價鼓勵書店在需求轉弱時仍持有庫存，這有助於競爭；但它也使需求低迷時無法降價出清，滯銷存貨因此累積，造成環境成本。我們建立含需求不確定性的動態寡占模型，以書店月度資料估計，探討有多少滯銷存貨源自價格承諾，以及押金退款稅能如何減少浪費。台灣公平交易委員會在《公平交易法》下將維持轉售價格視為原則違法，並曾處分固定書籍轉售價格的出版商，包括對天下遠見出版固定《賈伯斯傳》售價新臺幣599元處以新臺幣20萬元罰鍰，以及對龍騰限制書店參考書折扣處以新臺幣50萬元罰鍰。\footnote{公平交易委員會處分：公處字第101137號（2012年9月19日，天下遠見出版）；（89）公處字第207號（2000年11月30日，龍騰）。案例摘要：\url{https://www.apeccp.org.tw/htdocs/doc/Taipei/Case/dec2012/1136-1089case07.html}；\url{https://www.apeccp.org.tw/htdocs/doc/Taipei/Case/D0505200.htm}。}
\end{itemize}

\medskip
\noindent\textbf{未來研究計畫。}未來五年，我計畫將此議程發展為對「造成浪費的市場規則」更廣泛的分析，並與環境及循環經濟政策連結。產品退貨與滯銷存貨，是同一個市場失靈在銷售兩個階段的表現，兩者都受價格、費用與交易反悔權的形塑。我看到三個方向。

第一，我將共同研究搜尋、購買與退貨。搜尋與退貨是解決同一個媒合不確定性的兩種方式，一個在付款前、一個在付款後，因此互為替代。更好的搜尋降低買方帶進購買的剩餘不確定性，應可降低退貨率；但它也讓賣家可以收取更高價格，因為買方的外部選擇資訊更充分。哪個力量主導，決定平台的推薦系統是真的減少退貨，還是只是轉而提高價格。我希望估計這個聯合決策，還原搜尋強度與退貨率之間的替代關係。可能的成果是一套福利核算，將媒合品質的提升與價格上升分開，並檢驗本意是提升媒合品質的推薦系統，是否反而把媒合失敗的成本轉嫁到退貨運輸與處置上。在台灣，免費退貨是法定預設，\footnote{消費者保護法第19條，\url{https://law.moj.gov.tw/ENG/LawClass/LawAll.aspx?pcode=J0170001}。}搜尋與退貨之間的替代、以及平台推薦在其中的角色，正是消費者已經面對的市場樣貌。

第二，我將把退貨與滯銷存貨作為環境問題來研究，並評估減少浪費的政策工具。在服裝業，消費者一次訂購數個尺寸或顏色，留下一件、退回其餘，退回商品造成大量損失；在出版業，固定零售價產生被銷毀的滯銷書籍。兩者都是保護某一環節的規則在另一環節製造實體浪費，因為制定規則的一方不承擔處置成本。我計畫量化這些浪費，並比較將其內部化的工具，包括補貨費、退貨運費、押金退款稅與價格彈性規則。直觀上，每種工具都使退貨的私人成本接近社會成本，但它們在多大程度上同時抑制有價值的購買並不相同，因此比較將得出每犧牲一單位交易、哪種工具減少最多浪費的排序。台灣環境部已管制電商包裝：禁用PVC、規定紙箱回收紙含量90\%以上、塑膠包材再生料25\%以上，並要求大型業者達成包裝減重與重複使用目標；退貨則在這些流向之上，額外增加逆向物流與包裝廢棄物。\footnote{環境部《網際網路購物包裝限制使用對象及實施方式》，2023年7月1日施行。英文摘要：\url{https://www.moenv.gov.tw/en/news/press-releases/15507.html}。}隨著台灣延伸生產者責任與循環經濟規範趨嚴，這項研究可為電商退貨與滯銷存貨應如何計算、收費或抑制提供依據。

第三，我將研究退貨與退款規則如何形塑平台的雙邊。退貨政策是影響參與的槓桿，不只影響加入後的行為。寬鬆的政策吸引更多買方，包括更可能退貨的買方，這提高賣家面對的退貨率、降低賣家收益，於是部分賣家退出，留下的賣家重新定價。跨邊效果因此可能抵銷寬鬆政策原本要創造的部分需求。我計畫估計一方的組成如何影響另一方，並比較平台會選擇的退貨政策與監理機關會選擇的退貨政策。這將得出從跨邊網路效應到退貨運費與補貨費在買方、賣方與平台之間最適分配的對應。公平交易委員會對數位平台競爭的研究，已使平台為雙邊制定的規則在台灣受到檢視。\footnote{公平交易委員會委託研究，中華經濟研究院執行：《數位平臺產業現況、發展趨勢及競爭規範—以我國為中心》（民國111年），\url{https://www.ftc.gov.tw/upload/1f59ab5d-8cae-410a-a8b1-2080fdab977c.pdf}。}

\medskip
\noindent\textbf{對中華經濟研究院的貢獻。}我的領域是實證產業組織，而我的工具正是貴院已經在使用的工具。中華經濟研究院接受公平交易委員會委託進行競爭研究，包括網路銷售市場的實證評估，以及數位平台與競爭規範的研究。\footnote{公平交易委員會委託研究，均由中華經濟研究院執行：《網路銷售市場競爭評估之實證研究》（民國109年），\url{https://www.ftc.gov.tw/upload/885e4488-950e-4cf9-a774-6a224b8df58c.pdf}；《數位平臺產業現況、發展趨勢及競爭規範—以我國為中心》（民國111年），\url{https://www.ftc.gov.tw/upload/1f59ab5d-8cae-410a-a8b1-2080fdab977c.pdf}。}這些報告界定相關市場、評估多邊網路效應與市場力量，並檢視獨家交易、維持轉售價格、最惠客戶條款與平台平價等行為。這正是我研究的領域。我的訓練是處理豐富的微觀資料，估計需求、供給與廠商行為的結構模型，並用這些模型評估規則或價格的改變如何影響市場力量、消費者剩餘與福利。我將把這個能力帶進貴院的競爭與數位經濟研究。

我的數個計畫與此議程直接相關：轉嫁計畫探討購買可退時費用與租稅如何傳導到價格；滯銷存貨計畫估計維持轉售價格的競爭效果，亦即公平交易委員會指出的垂直限制之一；求職論文與搜尋計畫則處理數位平台研究中的定價、搜尋與平台設計。在這些工作中，我處理從原始交易或掃描資料到政策制定者可採納的政策模擬的完整路徑。

同一套工具也服務於貴院其他中心。對台灣產業與能源及綠色經濟相關中心，我在滯銷存貨與產品退貨上的研究，衡量定價規則與運費補貼造成的浪費，並評估可減少浪費的工具，如押金退款稅與補貨費，這與延伸生產者責任及循環經濟政策直接相關。對經濟預測，我的轉嫁研究顯示，忽略退貨邊際可能使成本衝擊的預測價格反應完全相反。我在台灣長大，能以中文與英文和國內資料及利害關係人合作。我在中華經濟研究院的目標，是把結構估計轉化為競爭、數位平台、租稅與環境政策上清楚、以證據為本的建議。

\end{document}
